\section{Secure Hash Algorithm}\label{sec:krypto-sha}

Zkratka SHA pochází z~anglického názvu \emph{Secure Hash Algorithm}. Neoznačuje jedinou funkci, ale několik rodin standardizovaných kryptografických hešovacích funkcí.

\begin{itemize}
    \item \textbf{SHA-0} byl zveřejněn v~roce 1993 jako původní Secure Hash Standard.
    \item \textbf{SHA-1} jej v~roce 1995 nahradil. Později byly nalezeny praktické kolize, a~proto se pro nové bezpečnostní použití nepoužívá.
    \item \textbf{SHA-2} je rodina funkcí zahrnující například SHA-224, SHA-256, SHA-384 a~SHA-512.
    \item \textbf{SHA-3} byl standardizován v~roce 2015 a~vychází z~algoritmu Keccak. Má odlišnou vnitřní konstrukci než SHA-2.
\end{itemize}

\begin{figure}[ht]
    \centering
    \begin{tikzpicture}[
    font=\normalsize,
    >=Latex,
    point/.style={circle, draw=blue!65!black, fill=blue!18,
                  minimum size=6mm, inner sep=0pt},
    label/.style={rounded corners=1mm, draw=black!45, fill=black!3,
                  minimum width=2.6cm, minimum height=10mm, align=center}
]
    \draw[-{Latex}, line width=1pt, draw=black!65] (-6,0) -- (6.4,0);
    \foreach \x/\year/\name/\desc in {
        -5/1993/SHA-0/původní standard,
        -2/1995/SHA-1/dnes nevhodný,
        1.5/2002/SHA-2/rodina funkcí,
        5/2015/SHA-3/Keccak
    }{
        \node[point] at (\x,0) {};
        \node[font=\bfseries] at (\x,.62) {\year};
        \node[label] at (\x,-1.05) {\name\\\desc};
    }
\end{tikzpicture}

    \caption{Zjednodušená časová osa rodin SHA.}
    \label{fig:krypto-sha-casova-osa}
\end{figure}

Pro nové aplikace se podle konkrétního standardu a~protokolu používají zejména funkce z~rodin SHA-2 a~SHA-3. Skutečnost, že obě rodiny zůstávají v~použití, není rozpor: mají jinou konstrukci a~mohou představovat nezávislé možnosti pro různé systémy.

\subsection{Sponge function}

SHA-3 používá konstrukci nazývanou \emph{sponge function}, doslova „houbová funkce“. Název připomíná dvě fáze: v~první funkce postupně „nasává“ bloky zprávy a~ve druhé ze svého vnitřního stavu „vymačkává“ výstupní bity.

Vnitřní stav má délku
\[
    b=r+c.
\]
Prvních $r$ bitů tvoří část \emph{rate}, která komunikuje se vstupem a~výstupem. Zbývajících $c$ bitů tvoří \emph{capacity}; nejsou přímo vydávány jako výstup a~významně ovlivňují bezpečnost konstrukce. Celý stav zpracovává veřejná kryptografická permutace
\[
    \mapping{f}{\set{0,1}^{b}}{\set{0,1}^{b}}.
\]

\subsection{Fáze absorbing}

Zprávu $m$ nejprve jednoznačně doplníme předepsaným paddingem a~rozdělíme na $r$-bitové bloky
\[
    p_1,p_2,\ldots,p_t.
\]
Začneme nulovým stavem $S_0=0^b$. Každý blok prodloužíme o~$c$ nulových bitů, zkombinujeme s~předchozím stavem operací XOR a~na výsledek použijeme permutaci $f$:
\[
    S_i
    =
    f\left(S_{i-1}\oplus(p_i0^c)\right),
    \qquad i=1,\ldots,t.
\]
Řetězec $p_i0^c$ má délku $r+c=b$, takže jej lze XORovat s~celým stavem. Vstupní blok ovlivní po aplikaci permutace obě části stavu.

\subsection{Fáze squeezing}

Po zpracování všech bloků čteme bity z~části rate výsledného stavu $S_t$. Potřebujeme-li více než $r$ výstupních bitů, znovu použijeme permutaci $f$ a~z~nového stavu odebereme další část rate. Proces opakujeme, dokud nemáme požadovanou délku otisku.

\begin{figure}[ht]
    \centering
    \begin{tikzpicture}[
    font=\normalsize,
    >=Latex,
    block/.style={rounded corners=1mm, draw=black!55, fill=black!4,
                  minimum width=2cm, minimum height=9mm, align=center},
    state/.style={rounded corners=1mm, draw=blue!65!black, fill=blue!11,
                  minimum width=1.8cm, minimum height=9mm, align=center},
    func/.style={rounded corners=1mm, draw=orange!75!black, fill=orange!18,
                 minimum width=10mm, minimum height=10mm},
    xor/.style={circle, draw=black!55, fill=white, minimum size=7mm, inner sep=0pt},
    flow/.style={-{Latex}, draw=black!70, line width=.75pt}
]
    \node[block] (m) at (-6.2,3) {zpráva \(m\)};
    \node[block] (pad) at (-3.15,3) {padding};
    \node[block, minimum width=3.2cm] (parts) at (.25,3)
        {bloky \(p_1,\ldots,p_t\)};
    \draw[flow] (m) -- (pad);
    \draw[flow] (pad) -- (parts);

    \node[state] (s0) at (-6.2,0) {\(S_0=0^b\)};
    \node[xor] (x1) at (-4.0,0) {\(\oplus\)};
    \node[func] (f1) at (-2.4,0) {\(f\)};
    \node[state] (s1) at (-.8,0) {\(S_1\)};
    \node at (.6,0) {\(\cdots\)};
    \node[xor] (xt) at (2.0,0) {\(\oplus\)};
    \node[func] (ft) at (3.6,0) {\(f\)};
    \node[state] (st) at (5.2,0) {\(S_t\)};

    \draw[flow] (s0) -- (x1);
    \draw[flow] (x1) -- (f1);
    \draw[flow] (f1) -- (s1);
    \draw[flow] (s1) -- (.25,0);
    \draw[flow] (.95,0) -- (xt);
    \draw[flow] (xt) -- (ft);
    \draw[flow] (ft) -- (st);

    \node[block] (p1) at (-4.0,1.25) {\(p_1 0^c\)};
    \node[block] (pt) at (2.0,1.25) {\(p_t 0^c\)};
    \draw[flow] (parts.south west) to[bend right=8] (p1.north);
    \draw[flow] (parts.south east) to[bend left=8] (pt.north);
    \draw[flow] (p1) -- (x1);
    \draw[flow] (pt) -- (xt);

    \node[block, draw=green!55!black, fill=green!12] (rate) at (5.2,-1.7)
        {čti část rate};
    \node[state, minimum width=2.6cm] (hash) at (5.2,-3.15)
        {otisk \(\mathcal{H}_n(m)\)};
    \draw[flow] (st) -- (rate);
    \draw[flow] (rate) -- (hash);

    \node[draw=blue!55!black, rounded corners=2mm, inner sep=4mm,
          fit=(s0)(p1)(st)(pt)] (absorbing) {};
    \node[font=\bfseries, text=blue!55!black, fill=white,
          inner xsep=2mm, inner ysep=.5mm]
        at ([xshift=1.8cm]absorbing.north west) {absorbing};

    \node[draw=green!55!black, rounded corners=2mm, inner sep=3mm,
          fit=(rate)(hash)] (squeezing) {};
    \node[font=\bfseries, text=green!45!black, fill=white,
          inner xsep=2mm, inner ysep=.5mm]
        at ([xshift=.9cm]squeezing.north west) {squeezing};
\end{tikzpicture}

    \caption{Fáze absorbing a~squeezing konstrukce sponge function.}
    \label{fig:krypto-sponge}
\end{figure}

\notebox{Permutace $f$ není tajný klíč. Je veřejnou a~pevně určenou součástí algoritmu. Bezpečnost má vycházet z~její konstrukce a~ze způsobu, jakým se celý vnitřní stav používá, nikoli z~utajení postupu.}

\subsection{SHA-3 jako konkrétní sponge function}

SHA-3 používá stav délky
\[
    b=1600.
\]
Pro variantu SHA3-$d$ s~$d$-bitovým otiskem se volí
\[
    c=2d,
    \qquad
    r=1600-c.
\]
Například SHA3-256 má
\[
    d=256,
    \qquad
    c=512,
    \qquad
    r=1088.
\]
Další běžné varianty jsou SHA3-224, SHA3-384 a~SHA3-512. Liší se délkou otisku, a~tím i~rozdělením stavu na rate a~capacity.

SHA-3 je navržen jako bezpečná kryptografická hešovací funkce, formální důkaz, že konkrétní SHA3-$d$ splňuje definici jednosměrné funkce z~\smartref[definice~][]{\showrefs}{def:krypto-owf}{předchozího výkladu}, však znám není. To není u~praktických kryptografických funkcí neobvyklé: bezpečnost se opírá o~rozsáhlou veřejnou analýzu, známé útoky a~přesně popsané bezpečnostní předpoklady.

\importantbox{Délka otisku a~vnitřní konstrukce řeší různé otázky. Delší otisk zvětšuje prostor možných hodnot, ale bezpečnost stále závisí také na tom, zda ve funkci neexistuje strukturální zkratka rychlejší než obecný útok hrubou silou.}
